\section{Discussion}\label{sec:discussion}

Dynamic-DeepHit has the highest recorded mean discrimination across all three scenarios, followed by Hazard Transformer. The completed twenty-feature results replace the pilot findings. However, feature count varies with cohort composition, missingness representation, and access to patient histories; these experiments cannot isolate the benefit of additional measurements or attribute performance differences to architecture alone.

Native calibration is less consistent than discrimination, and the common transformed-score IBS does not measure native probability calibration. The decision-curve comparisons require harmonized patient-level baselines before any referral benefit can be inferred.

\subsection{Limitations}\label{sec:limitations}
Five runs on a fixed partition describe fitting/tuning variability, not patient-sampling uncertainty; there are no patient-bootstrap confidence intervals or external validation. The event-rich hospital cohort limits generalizability. IPCW uses raw training rows, and Logistic Hazard uses the nominal test set in tuning/checkpoint selection. These evaluation choices require validation in a rerun with separate tuning data and patient-level censoring estimates.

Models also differ in access to longitudinal histories. Final-observation scoring, whole-history uACR matching, and bidirectional chemistry windows do not establish prospective performance at a common decision time. The uACR match attaches a value recorded after the anchor creatinine for 40\% of rows (median absolute separation 264 days, p95 5.4 years), whereas the $\pm$24h chemistry matches are simultaneous with the anchor for 93--99\% of rows. For label-positive patients the median time from the first extracted record to the labelled event is 1.2 days (four-feature), 0.4 days (eight-feature) and 104 days (twenty-feature), and 18--35\% contribute a single row, so a two-year horizon measured from that origin is a follow-up window rather than a prediction lead time. Cox remains below chance across scenarios, eight-feature DeepSurv averages below 0.5, and Hazard Transformer has constant two-year probabilities; shared implementation effects remain possible. The labelled outcome is an acute/unspecified kidney-failure code, not the dialysis-or-transplant endpoint KFRE predicts (Section~\ref{sec:outcome}), so KFRE-versus-model calibration differences combine an endpoint difference, a population difference, and the fitted/externally-specified asymmetry, which this benchmark does not separate. Death before ESRD is not modeled as a competing event, demographic subgroup performance is untested, and twenty-feature DDH/Transformer importance outputs are failed-probe fallbacks. Future work should address these issues before clinical utility is assessed.
